Reinforcement Learning with Verifiable Rewards (RLVR) has recently demonstrated notable success in enhancing the reasoning performance of large language models (LLMs), particularly in mathematics and programming tasks. 
It is widely believed that, similar to how traditional RL helps agents to explore and learn new strategies, RLVR enables LLMs to continuously self-improve, thus acquiring novel reasoning abilities that exceed the capacity of the corresponding base models. 
In this study, we take a critical look at \textit{the current state of RLVR} by systematically probing the reasoning capability boundaries of RLVR-trained LLMs across various model families, RL algorithms, and math/coding/visual reasoning benchmarks, using pass@\textit{k} at large \textit{k} values as the evaluation metric.
While RLVR improves sampling efficiency towards correct paths, we surprisingly find that current training \emph{rarely} elicit fundamentally new reasoning patterns.
We observe that while RLVR-trained models outperform their base models at smaller values of $k$ (\eg, $k$=1), base models achieve higher pass@$k$ score when $k$ is large.
Moreover, we observe that the reasoning capability boundary of LLMs often narrows as RLVR training progresses.
Further coverage and perplexity analysis shows that the reasoning paths generated by RLVR models are already included in the base models' sampling distribution, suggesting that their reasoning abilities originate from and are \textit{bounded} by the base model. 
From this perspective, treating the base model as an upper bound, our quantitative analysis shows that six popular RLVR algorithms perform similarly and remain far from optimal in fully leveraging the potential of the base model.
In contrast, we find that distillation can introduce new reasoning patterns from the teacher and genuinely expand the model's reasoning capabilities.
Taken together, our findings suggest that current RLVR methods have not fully realized the potential of RL to elicit genuinely novel reasoning abilities in LLMs. This underscores the need for improved RL paradigms, such as effective exploration mechanism, more deliberate and large-scale data curation, fine-grained process signal, and multi-turn agent interaction, to unlock this potential.

% Surprisingly, we found that current RLVR training does \emph{not}, in fact, elicit fundamentally new reasoning patterns. We observed that while RL-trained models outperform their base models at smaller values of $k$ (\eg, $k$=1), base models can achieve a comparable or even higher pass@$k$ score compared to their RL counterparts at large $k$ values.
% Further analysis shows that the reasoning paths generated by RL-trained models are already included in the base models' sampling distribution, suggesting that most reasoning abilities manifested in RL-trained models are already obtained by base models. 
% RL training boosts the performance by biasing the model's output distribution toward paths that are more likely to yield rewards, therefore sampling correct responses more efficiently. But this also limits their exploration capacity, resulting in a narrower reasoning capability boundary compared to base models.
% Surprisingly, we found that current RLVR training does \emph{not}, in fact, elicit fundamentally new reasoning patterns. 

